\documentclass[10pt,a4paper]{amsart}
\usepackage[margin=19mm]{geometry}
\usepackage{amsmath,amssymb,mathtools}
\usepackage[scr=rsfs,cal=euler]{mathalfa}
\usepackage{xfrac}
\usepackage[unicode,hidelinks,bookmarks=false]{hyperref}
\newcommand{\ie}{i.e.\ }
\newcommand{\cf}{cf.\ }
\newcommand{\resp}{respectively\ }
\newcommand{\Subsection}{\subsection}
\providecommand{\nobreakdash}{\nobreak\hbox{-}}
\setlength{\emergencystretch}{2em}
\multlinegap0pt
\usepackage{bm}
\usepackage{bbm}
\usepackage{dsfont}
\usepackage{xspace}
\usepackage{cancel}
\usepackage{mathtools}
\usepackage{amsthm}
\makeatletter
\@ifundefined{theorem}{\newtheorem{theorem}{Theorem}}{}
\@ifundefined{lemma}{\newtheorem{lemma}[theorem]{Lemma}}{}
\makeatother
\DeclareMathOperator*{\Res}{Res}
\renewcommand{\Im}{\operatorname{Im}}
\renewcommand{\Re}{\operatorname{Re}}
\title[Self-Correlations of Hurwitz Class Numbers]{Self-Correlations of Hurwitz Class Numbers}
\author{Alexander Walker}
\date{Source: arXiv:2402.01455v1 (CC BY 4.0)}
\begin{document}
\maketitle

\noindent\emph{Source and licence.} Self-Correlations of Hurwitz Class Numbers. Version arXiv:2402.01455v1 (CC BY 4.0), distributed under the Creative Commons Attribution 4.0 International licence, as recorded in the arXiv licence metadata for this exact version and shown on the abstract page https://arxiv.org/abs/2402.01455v1. Full rights notice: licenses/arxiv-2402.01455-CC-BY-4.0.txt.\par\smallskip

\noindent\textit{Editorial scope.} Counted unit: the Lemma whose source label is \texttt{lem:phi-cross-behavior} in the article (source0.tex lines 1261--1279, source label \texttt{lem:phi-cross-behavior}), delivered together with the article's own complete original proof of it (source0.tex lines 1281--1311, one \texttt{proof} environment of 31 lines). No mathematics is written, repaired or re-derived: statement and proof are the article's own text line for line. Required context retained uncounted: eq:incomplete-gamma-integral-representation (lines 935--939); eq:phi-definitions (lines 1119--1130); eq:phi-plus-closed-form (lines 1188--1192); lem:phi-plus-residue-bound (lines 1235--1241). Every definition, display and label of those lines is retained unchanged. Omitted from this package and not redistributed: 1--934, 940--1118, 1131--1187, 1193--1234, 1242--1260, 1280--1280, 1312--1838. Nothing inside a retained excerpt is omitted or altered: every retained body is delivered exactly as the article's lines are written, so no omission receipt of any kind accompanies this package. Citations: the retained excerpts carry no citation command at all, so no bibliography entry of the article is transcribed and no reference is redistributed. Reference adaptation: none was needed and none was made; every reference inside the delivered unit resolves inside it, so no unresolved reference or citation is produced. Printed slips kept literal: no printed word, symbol, hypothesis or number of the counted statement or of its proof was altered. Layout and class: the article's own class is \texttt{amsart} with packages including inputenc, amsmath, amsthm, amssymb, mathrsfs, graphics, hyperref, xcolor, soul, mathtools, natbib, verbatim, longtable; the wrapper is the lane's pinned \texttt{amsart} editorial wrapper at 10pt on A4, which restates the theorem-like environments the retained text needs and adds the article's own macro definitions that the retained text needs (\texttt{\textbackslash Im}, \texttt{\textbackslash Re}, \texttt{\textbackslash Res}), with extra wrapper packages (bm, bbm, dsfont, xspace, cancel, mathtools). The delivered numbering is the unit's own. Assets: no retained span contains a figure environment, a graphics-inclusion command, a PDF-page inclusion, a table or a TikZ picture; no image file is copied into this package, the package declares no asset dependency and no image file is redistributed with it. Licence: version arXiv:2402.01455v1 of this article is distributed under the Creative Commons Attribution 4.0 International licence, as recorded in the arXiv licence metadata for this exact version and shown on the abstract page https://arxiv.org/abs/2402.01455v1. A full rights notice is delivered at licenses/arxiv-2402.01455-CC-BY-4.0.txt. Characters: every retained excerpt is pure ASCII, so the delivered engine, XeLaTeX with its default Latin Modern text font, reports no missing character.
\begin{align} \label{eq:incomplete-gamma-integral-representation}
	\Gamma(1-k, y) e^y =
		- \frac{y^{-k}}{\Gamma(k)} \cdot \frac{\pi}{2\pi i}
		\int_C \frac{\Gamma(w+k) y^{-w}}{\sin(\pi w)} \,dw,
\end{align}

\begin{align}
	\varphi_j^+(m,n,s) 
		&:= \int_0^\infty \! y^{s+k-\frac{1}{2}} e^{-2\pi (2n+m) y}
				K_{it_j}(2\pi m y) \frac{dy}{y}, \\ \label{eq:phi-definitions}
	\varphi_j^-(m,n,s)
		&:= \int_0^\infty \! y^{s+k-\frac{1}{2}} e^{2\pi (2n+m) y} \Gamma(1-k, 4\pi n y) \\[-5pt]
			& \hspace{30 mm} \times
			\Gamma(1-k, 4\pi(n+m) y)  K_{it_j}(2\pi m y) \frac{dy}{y}, \\
	\varphi_j^\times(m,n,s)
		&:= \int_0^\infty \! y^{s+k-\frac{1}{2}} e^{2\pi(m-2n)y} 
			\Gamma(1-k, 4\pi (m-n)y) K_{it_j}(2\pi m y) \frac{dy}{y}.
\end{align}

\begin{align} \label{eq:phi-plus-closed-form}
	\varphi_j^+(m,0,s)
		= \frac{\sqrt{\pi} \, \Gamma(p+it_j)\Gamma(p-it_j)}
					{(4\pi m )^p  \Gamma(p+\frac{1}{2})},
\end{align}

\begin{lemma} \label{lem:phi-plus-residue-bound}
Fix $t_j \in \mathbb{R}$. For each integer $r \geq 0$, we have
\[
	\Res_{s =\frac{1}{2}-k \pm it_j - r} \varphi_j^+(m,n,s)
	\ll_r (n+m)^r \vert t_j \vert^{-\frac{1}{2}} e^{-\frac{\pi}{2} \vert t_j \vert}.
\]
\end{lemma}

\begin{lemma} \label{lem:phi-cross-behavior}
The function $\varphi_j^\times(m,n,s)$ defined in~\eqref{eq:phi-definitions} admits meromorphic continuation to $s \in \mathbb{C}$, with poles at $s = -\frac{1}{2} \pm it_j -r$ and $s = \frac{1}{2} - k \pm it_j - r$, for $r \in \mathbb{Z}_{\geq 0}$. If $t_j \in \mathbb{R}$ and $\Re s > 0$ away from poles, we have
\begin{align} \label{eq:phi-cross-upper-bound}
	\varphi_j^\times(m,n,s)
	&\ll  \frac{1}{(m-n)^k\sqrt{m}}
	\Big( \frac{\vert s - it_j \vert \cdot \vert s + it_j\vert}{m \vert s \vert}\Big)^{\Re s -1}
	\vert s \vert^{-\frac{1}{2}} \\
	& \qquad
	\times \! \Big( 
		1 + \Big( 
					\frac{m\vert s \vert}{(m-n)\vert s - it_j \vert \cdot \vert s + it_j\vert}
					\Big)^{\vert k \vert +\epsilon} \Big)
	e^{-\pi \vert t_j \vert + \frac{\pi}{2} \vert \Im s \vert}  \\
	& \quad 
	+ \delta_{[\Re s < \frac{1}{2} - k]} \big(m+n+ \vert s \vert + \vert t_j \vert\big)^A
	e^{-2\pi \vert t_j \vert + \frac{3\pi}{2} \vert \Im s \vert}
\end{align}
for all $\epsilon > 0$ and for some $A>0$ depending only on $k$.
\end{lemma}

\begin{proof}
The integral representation~\eqref{eq:incomplete-gamma-integral-representation} implies that
\begin{align}
	\varphi_j^\times(m,n,s)
		&=
			\frac{-\pi}{\Gamma(k)} 
			\cdot  \frac{1}{2\pi i} \int_C \frac{\Gamma(w+k)}{(4\pi(m-n))^{w+k} \sin(\pi w)} \\
		& \qquad\qquad \qquad
	\times\Big( \int_0^\infty \! y^{s-\frac{1}{2}-w} e^{-2\pi m y} K_{it_j}(2\pi m y) \frac{dy}{y} \Big) dw \\
		\label{eq:phi-cross-integral-representation}
		&= 
			\frac{-\pi}{\Gamma(k)} \cdot  \frac{1}{2\pi i} 
			\int_C \frac{\Gamma(w+k)\varphi_j^+(m,0, s-k-w)}{(4\pi(m-n))^{w+k} \sin(\pi w)} dw,
\end{align}
where $C$ is a contour separating the poles of $\Gamma(w+k)$ from those at $w=0,1,\ldots$ arising from $1/\sin(\pi w)$. To begin, we require $\Re s > \frac{1}{2} + \max \{\Re w: w \in C\}$. To consider general $s$, we shift the contour $C$ left, passing poles from $\Gamma(w+k)$ and $1/\sin(\pi w)$ and extracting residues involving $\varphi_j^+(m,0,s)$ at shifted arguments. By~\eqref{eq:phi-plus-closed-form}, these residues contribute poles at the poles of $\Gamma(s+k-\frac{1}{2} \pm it_j)$ and $\Gamma(s + \frac{1}{2} \pm it_j)$.

To produce growth estimates, we then shift $C$ rightwards, to the contour $\Re w = \vert k \vert + \epsilon$. This extracts a sum of residues equal to
\begin{align*}
	& \sum_{q = 0}^{\vert k \vert - \frac{1}{2}} 
			\frac{(-1)^q\Gamma(q+k)}{\Gamma(k) (4\pi(m-n))^{k+q}} \cdot \varphi_j^+(m,0,s-k-q) \\[-.38em]
	& \qquad + \!\!\!\!
		\sum_{r=0}^{\lfloor \vert k \vert +\epsilon - \Re s + \frac{1}{2} \rfloor} \!\!\!\!
		\sum_{\pm}
			\Res_{w = s-\frac{1}{2} \pm it_j +r}
				\frac{\pi \Gamma(w+k)\varphi_j^+(m,0, s-k-w)}{(4\pi(m-n))^{w+k} \sin(\pi w) \Gamma(k)}.
\end{align*}
Stirling's approximation and Lemma~\ref{lem:phi-plus-residue-bound} show that the exponential decay in the residues in the second line is $e^{-\frac{3\pi}{2} \vert \Im s \pm it_j \vert - \frac{\pi}{2} \vert t_j \vert} \ll e^{-2\pi \vert t_j \vert + \frac{3\pi}{2} \vert \Im s \vert}$, while the worst polynomial growth is $O((m+n+ \vert s \vert + \vert t_j \vert)^A)$ for some $A> 0$ depending linearly on $\vert k \vert$ and $\Re s $. Since $\Re s \in (0, \vert k \vert + \frac{1}{2})$ when these terms appear, we may take the constant $A$ to depend on $k$ alone.


Exponential decay in $\vert \Im w \vert$ within the integrand bounds the shifted contour integral to at most a constant multiple of the integrand near $\vert k \vert + \epsilon$. Stirling's approximation and~\eqref{eq:phi-plus-closed-form} then complete the proof of~\eqref{eq:phi-cross-upper-bound}.
\end{proof}

\end{document}
